\documentclass[11pt]{article}

\usepackage[utf8]{inputenc}
\usepackage[T1]{fontenc}
\usepackage{lmodern}
\usepackage{geometry}
\usepackage{amsmath,amssymb,amsfonts,amsthm,mathtools}
\usepackage{bm}
\usepackage{enumitem}
\usepackage{microtype}
\usepackage{hyperref}
\usepackage{titlesec}
\usepackage{booktabs}
\usepackage{array}
\usepackage{setspace}
\usepackage{mathrsfs}
\usepackage{physics}

\geometry{margin=1in}
\hypersetup{colorlinks=true, linkcolor=black, urlcolor=blue, citecolor=black}
\setlist[enumerate]{topsep=0.4em,itemsep=0.35em}
\setlist[itemize]{topsep=0.3em,itemsep=0.25em}
\titleformat{\section}{\large\bfseries}{\thesection.}{0.5em}{}
\titleformat{\subsection}{\normalsize\bfseries}{\thesubsection.}{0.5em}{}
\setstretch{1.06}

\newcommand{\Aether}{\AE{}ther}
\newcommand{\Aflow}{\AE{}ther-flow}
\newcommand{\TheoryName}{\Aflow{} Interpretation of Relativity}
\newcommand{\TheoryProgram}{\Aether{} / \Aflow{} framework}

\newcommand{\CompletedMetric}{g^{\sharp}}
\newcommand{\MatterSpace}{\Psi}

\newcommand{\Car}{\mathrm{car}}
\newcommand{\Cur}{\mathrm{cur}}
\newcommand{\Hom}{\mathrm{hom}}

\newcommand{\ForSpace}[1]{\mathcal I_{#1}^{\mathrm{for}}}
\newcommand{\PrimShell}{\mathcal S_{\mathrm{prim}}}
\newcommand{\Reduction}[1]{\mathcal R_{#1}}
\newcommand{\CurrentCarrier}[1]{\mathfrak C_{#1}^{\mathrm{cur}}}
\newcommand{\EqCarrier}[1]{\mathfrak C_{#1}^{\mathrm{eq}}}
\newcommand{\MetricOut}{\mathscr G_{\mathrm{out}}}
\newcommand{\MatterOut}{\mathscr T_{\mathrm{out}}}

\newcommand{\SubTarget}[1]{E_{#1}^{\mathrm{sub}}}
\newcommand{\EqnTarget}[1]{Q_{#1}^{\mathrm{eqn}}}
\newcommand{\HiddenSpace}[1]{\mathcal H_{#1}^{\mathrm{hid}}}
\newcommand{\HiddenBody}[1]{D_{#1}^{\mathrm{hid}}}

\newcommand{\SeedState}[1]{\mathcal S_{#1}^{\mathrm{seed}}}
\newcommand{\SeedBody}[1]{\mathcal B_{#1}^{\mathrm{seed}}}
\newcommand{\SeedShell}[1]{\mathcal Z_{#1}^{\mathrm{seed}}}
\newcommand{\SeedSupportShell}[1]{\mathcal U_{#1}^{\mathrm{seed,sup}}}
\newcommand{\SeedZeroCore}[1]{\mathcal Z_{#1}^{0,\mathrm{seed}}}
\newcommand{\SeedSplit}[1]{\Phi_{#1}^{\mathrm{seed}}}
\newcommand{\SeedCharge}[1]{\mathcal Q_{#1}^{\mathrm{seed}}}
\newcommand{\SeedEmbed}[1]{\zeta_{#1}^{\mathrm{seed}}}
\newcommand{\SeedKernel}[1]{\widetilde K_{#1}^{0,\mathrm{seed}}}
\newcommand{\SeedCoercive}[1]{P_{#1}^{\mathrm{seed,co}}}

\newcommand{\GermNucleus}{\mathfrak G^{\mathrm{sub,germ,zstn}}}
\newcommand{\GermState}[1]{\mathcal G_{#1}^{\mathrm{germ}}}
\newcommand{\GermBody}[1]{\mathcal B_{#1}^{\mathrm{germ}}}
\newcommand{\GermProj}[1]{\Pi_{#1}^{\mathrm{germ}}}
\newcommand{\GermSeedSec}[1]{\sigma_{#1}^{\mathrm{germ\leftarrow seed}}}
\newcommand{\GermSection}[1]{\Gamma_{#1}^{\mathrm{germ}}}
\newcommand{\GermShell}[1]{\mathcal Z_{#1}^{\mathrm{germ}}}

\newcommand{\ProtoTransportCore}{\mathfrak P^{\mathrm{sub,proto,xtc}}}
\newcommand{\ProtoSectionCore}{\mathfrak P^{\mathrm{sub,proto,sec}}}
\newcommand{\ProtoState}[1]{\mathcal P_{#1}^{\mathrm{proto}}}
\newcommand{\ProtoBody}[1]{\mathcal B_{#1}^{\mathrm{proto}}}
\newcommand{\ProtoProj}[1]{\Pi_{#1}^{\mathrm{proto}}}
\newcommand{\ProtoSection}[1]{\Gamma_{#1}^{\mathrm{proto}}}
\newcommand{\ProtoShell}[1]{\mathcal Z_{#1}^{\mathrm{proto}}}
\newcommand{\ProtoTransportSec}[1]{\mathfrak s_{#1}^{\mathrm{proto,sec}}}

\newcommand{\ProtoSectionPackage}{\mathfrak Q^{\mathrm{sub,proto,secgen}}}
\newcommand{\ProtoSectionState}[1]{\mathcal Q_{#1}^{\mathrm{psec}}}
\newcommand{\ProtoSectionBody}[1]{\mathcal B_{#1}^{\mathrm{psec}}}
\newcommand{\ProtoSectionProj}[1]{\Pi_{#1}^{\mathrm{psec}}}
\newcommand{\ProtoSectionFiber}[2]{\mathcal F_{#1}^{\mathrm{psec}}(#2)}
\newcommand{\ProtoSectionFunctional}[1]{\mathscr E_{#1}^{\mathrm{psec}}}
\newcommand{\ProtoSectionShell}[1]{\mathcal Z_{#1}^{\mathrm{psec}}}
\newcommand{\ProtoSectionMin}[1]{\mathfrak m_{#1}^{\mathrm{psec}}}
\newcommand{\InducedTransportSec}[1]{\widehat{\mathfrak s}_{#1}^{\mathrm{proto,sec}}}

\newcommand{\PreProtoSectionPackage}{\mathfrak R^{\mathrm{sub,proto,presec}}}
\newcommand{\PreProtoSectionState}[1]{\mathcal R_{#1}^{\mathrm{ppsec}}}
\newcommand{\PreProtoSectionBody}[1]{\mathcal B_{#1}^{\mathrm{ppsec}}}
\newcommand{\PreProtoSectionProj}[1]{\Pi_{#1}^{\mathrm{ppsec}}}
\newcommand{\PreProtoSectionFiber}[2]{\mathcal F_{#1}^{\mathrm{ppsec}}(#2)}
\newcommand{\PreProtoSectionProtoMin}[1]{\mathfrak n_{#1}^{\mathrm{ppsec}}}
\newcommand{\PreProtoSectionFunctional}[1]{\mathscr F_{#1}^{\mathrm{ppsec}}}
\newcommand{\PreProtoSectionShell}[1]{\mathcal Z_{#1}^{\mathrm{ppsec}}}

\newtheorem{definition}{Definition}
\newtheorem{proposition}{Proposition}
\newtheorem{remark}{Remark}
\newtheorem{corollary}{Corollary}
\newtheorem{theorem}{Theorem}

\title{\textbf{The \TheoryName{}}\\[0.4em]
\large Primitive-Reservoir Observer-Reduction\\
\large Exact-Shell-Transport-Section-Generating\\
\large Substrate Proto-Section Dynamics\\
\large from Proto-Section-Generating Substrate Pre-Proto-Section Dynamics}
\author{}
\date{}

\begin{document}

\maketitle

\begin{abstract}
The current primitive-reservoir observer-reduction frontier already moved the live burden below the germ-anchored exact-shell transport-section core itself. The active-primary theorem showed that the current section core is canonically generated by a deeper exact-shell-transport-section-generating substrate proto-section dynamics package \(\ProtoSectionPackage\). The accompanying routing layer then fixed the next honest move: either derive or justify that current proto-section package from still more primitive substrate structure, or prove a genuinely smaller theorem-level collapse, sharper necessity result, or sharper uniqueness result strictly inside it.

The present note takes the preferred first route. For each sector it introduces \emph{proto-section-generating substrate pre-proto-section dynamics}: a compact convex pre-proto-section body over the recorded proto-section body, an affine projection to that fixed body, a nonnegative smooth pre-proto-section functional with unique fiberwise minimizers, and an exact pre-proto-section shell projecting diffeomorphically to the recorded exact proto-section shell. The current proto-section functional is then no longer an unexplained primitive datum. It is the effective fiberwise minimum value induced by that deeper pre-proto-section law, and the current proto-section shell is the projected exact shell of the deeper package.

The benchmark boundary remains fixed throughout. The one operative metric is still exactly \(\CompletedMetric\), the ordinary matter route is unchanged, there is no second operative metric, the low-energy matter law is not enlarged, and no stronger promotion language is licensed. The positive result is therefore narrow but benchmark facing: the exact-shell-transport-section-generating substrate proto-section dynamics are no longer left as the primitive stopping point on the active line. The remaining honest burden moves one level deeper again, to the proto-section-generating substrate pre-proto-section dynamics themselves or to a sharper collapse, necessity theorem, or uniqueness theorem strictly inside that deeper package.
\end{abstract}

\section{Introduction}

The active derivational program of \TheoryName{} has again reached a sharper burden than another same-output relay. The current active-primary theorem already showed that the germ-anchored exact-shell transport-section core \(\ProtoSectionCore\) is not primitive: it is the canonical projected exact-shell minimizer selected by a deeper exact-shell-transport-section-generating substrate proto-section dynamics package \(\ProtoSectionPackage\). The routing consequence is immediate. Another deeper-origin manuscript that merely preserves that same proto-section package would not count as new front-line progress. The next honest benchmark-facing move must therefore either derive or justify \(\ProtoSectionPackage\) itself from still more primitive substrate structure, or prove a genuinely smaller theorem-level collapse, sharper necessity result, or sharper uniqueness result strictly inside that package.

The present note takes the direct derivation route. Its positive claim is not that completed substrate microphysics has now been identified, and it is not that the benchmark exact-relativistic package has changed. Its claim is narrower and benchmark facing. The current proto-section package is shown to arise canonically from a more primitive pre-proto-section layer whose projected exact shell and fiberwise minimum-value law already contain the benchmark-relevant structure of \(\ProtoSectionPackage\). In that precise sense the current proto-section dynamics are no longer the primitive place where the active line has to stop.

Two features matter here. First, the theorem targets the current proto-section package itself rather than re-spending the already settled section-core handoff, so it does not count progress merely by reproducing the same generated section core from yet another relay. Second, the deeper pre-proto-section package is not merely a renaming of the current proto-section package. It adds an explicit pre-proto-section state space, a pre-proto-section body, a pre-proto-to-proto projection, and a deeper fiberwise minimizing law whose effective energy reproduces the recorded proto-section functional on the benchmark-relevant body.

\paragraph{Framework and claim boundary.}
\TheoryProgram{} denotes the broader \Aether{} / \Aflow{} framework built on the claims that the \Aether{} is the underlying four-dimensional substrate of reality, the \Aflow{} is its intrinsic ordered motion, observed three-dimensional space is the local experiential slice of that deeper substrate, S-time is the experienced order of change arising from matter, light, and the \Aflow{}, and observed expansion is the three-dimensional appearance of deeper four-dimensional ordered motion. Within that framework, \TheoryName{} denotes the exact-closure theory in which the effective gravitational dynamics are adopted to be exactly Einsteinian with universal matter coupling. In the active sequence, \emph{adoption} means use of that established relativistic dynamics without claiming substrate derivation, while \emph{derivation} is reserved for a first-principles recovery from explicit substrate variables.

The present note stays strictly inside that boundary. It does not alter the benchmark exact-closure package. It does not introduce a second operative metric or a new low-energy matter law. It only proves that the current exact-shell-transport-section-generating substrate proto-section dynamics can itself be generated from a more primitive pre-proto-section-level substrate dynamics.

\section{Fixed Inputs}

\begin{definition}[Recorded primitive continuation data]
\label{def:fixed}
Fix the following continuation data from the active primitive-reservoir observer-reduction line.
\begin{enumerate}
    \item For each sector label \(\sigma\in\{\Car,\Cur,\Hom\}\), a primitive forcing shell
    \[
    \ForSpace{\sigma},
    \]
    a visible reduction map
    \[
    \Reduction{\sigma}:\ForSpace{\sigma}\to X_{\sigma},
    \]
    and shell-level current and equilibrium carriers
    \[
    \CurrentCarrier{\sigma}:\ForSpace{\sigma}\to Z_{\sigma}^{\mathrm{cur}},
    \qquad
    \EqCarrier{\sigma}:\ForSpace{\sigma}\to Z_{\sigma}^{\mathrm{eq}}.
    \]
    \item The product shell
    \[
    \PrimShell
    \equiv
    \ForSpace{\Car}\times \ForSpace{\Cur}\times \ForSpace{\Hom}.
    \]
    \item The recorded metric and ordinary-matter outputs
    \[
    \MetricOut:\PrimShell\to\{\text{observer metrics}\},
    \qquad
    \MatterOut:\PrimShell\times\MatterSpace\to\{\text{observer stress tensors}\},
    \]
    satisfying
    \begin{equation}
    \MetricOut(\mathbf w)=\CompletedMetric,
    \qquad
    \MatterOut(\mathbf w,\psi)
    =
    T_{\mu\nu}^{\mathrm{matter}}[\CompletedMetric,\psi]
    \label{eq:fixedoutputs}
    \end{equation}
    for every \(\mathbf w\in\PrimShell\) and every matter configuration \(\psi\in\MatterSpace\).
\end{enumerate}
\end{definition}

\begin{definition}[Recorded downstream seed package]
\label{def:seed}
Fix \(\sigma\in\{\Car,\Cur,\Hom\}\). The active line already fixes the downstream seed package at current benchmark-facing scope:
\begin{enumerate}
    \item a finite-dimensional seed state space \(\SeedState{\sigma}\), a compact convex seed body \(\SeedBody{\sigma}\subset\SeedState{\sigma}\), the exact seed shell \(\SeedShell{\sigma}\subset\SeedBody{\sigma}\), and the exact seed support shell \(\SeedSupportShell{\sigma}\subset\SeedShell{\sigma}\);
    \item a seed zero core \(\SeedZeroCore{\sigma}\subset\SeedSupportShell{\sigma}\), the same hidden fiber space \(\HiddenSpace{\sigma}\) with compact hidden body \(\HiddenBody{\sigma}\subset\HiddenSpace{\sigma}\), and the fixed seed split
    \[
    \SeedSplit{\sigma}:
    \SeedZeroCore{\sigma}\times\HiddenBody{\sigma}
    \xrightarrow{\cong}
    \SeedSupportShell{\sigma};
    \]
    \item the seed hidden charge
    \[
    \SeedCharge{\sigma}:\HiddenSpace{\sigma}\to\EqnTarget{\sigma},
    \]
    a smooth observer-compatible exact seed zero-response realization
    \[
    \SeedEmbed{\sigma}:\ForSpace{\sigma}\hookrightarrow\SeedZeroCore{\sigma},
    \]
    and the fixed-data solvability / coercive package recorded through \(\SeedKernel{\sigma}\) and \(\SeedCoercive{\sigma}\);
    \item benchmark faithfulness, meaning preservation of the benchmark outputs \eqref{eq:fixedoutputs}, no second operative metric, no enlarged low-energy matter law, no change to the ordinary matter route, and no premature promotion from adoption to completed derivation.
\end{enumerate}
\end{definition}

\begin{definition}[Recorded current minimal germ zero-shell transport nucleus]
\label{def:nucleus}
Fix \(\sigma\in\{\Car,\Cur,\Hom\}\). The recorded current minimal germ zero-shell transport nucleus \(\GermNucleus\) for sector \(\sigma\) consists of:
\begin{enumerate}
    \item the germ state space \(\GermState{\sigma}\), compact convex germ body \(\GermBody{\sigma}\subset\GermState{\sigma}\), projection
    \[
    \GermProj{\sigma}:\GermState{\sigma}\to\SeedState{\sigma},
    \]
    and seed section
    \[
    \GermSeedSec{\sigma}:\SeedBody{\sigma}\to\GermBody{\sigma}
    \]
    satisfying
    \[
    \GermProj{\sigma}\bigl(\GermBody{\sigma}\bigr)=\SeedBody{\sigma},
    \qquad
    \GermProj{\sigma}\circ\GermSeedSec{\sigma}
    =
    \mathrm{id}_{\SeedBody{\sigma}};
    \]
    \item the restriction section
    \[
    \GermSection{\sigma}:\GermState{\sigma}\to\SubTarget{\sigma},
    \]
    and the exact germ shell
    \[
    \GermShell{\sigma}
    \equiv
    \bigl(\GermSection{\sigma}\bigr)^{-1}(0)\cap\GermBody{\sigma},
    \]
    with
    \[
    \GermProj{\sigma}\big|_{\GermShell{\sigma}}
    :
    \GermShell{\sigma}
    \xrightarrow{\cong}
    \SeedShell{\sigma};
    \]
    \item benchmark faithfulness, meaning preservation of the benchmark outputs \eqref{eq:fixedoutputs}, no second operative metric, no enlarged low-energy matter law, no change to the ordinary matter route, and no premature promotion from adoption to completed derivation.
\end{enumerate}
By the already-recorded current frontier theorem, the support shell, zero core, hidden split, hidden charge, exact zero-response realization, and fixed-data solvability / coercive package are then canonically induced downstream from this nucleus together with the fixed seed package of Definition \ref{def:seed}, rather than taken as additional independent benchmark-facing data.
\end{definition}

\begin{definition}[Recorded current germ-anchored exact-shell transport core]
\label{def:core}
Fix \(\sigma\in\{\Car,\Cur,\Hom\}\). The recorded current germ-anchored exact-shell transport core \(\ProtoTransportCore\) for sector \(\sigma\) consists of:
\begin{enumerate}
    \item the proto-germ state space \(\ProtoState{\sigma}\), compact convex proto-germ body \(\ProtoBody{\sigma}\subset\ProtoState{\sigma}\), and affine surjection
    \[
    \ProtoProj{\sigma}:\ProtoState{\sigma}\to\GermState{\sigma}
    \]
    satisfying
    \[
    \ProtoProj{\sigma}\bigl(\ProtoBody{\sigma}\bigr)=\GermBody{\sigma};
    \]
    \item the restriction section
    \[
    \ProtoSection{\sigma}:\ProtoState{\sigma}\to\SubTarget{\sigma},
    \]
    the exact proto-germ shell
    \[
    \ProtoShell{\sigma}
    \equiv
    \bigl(\ProtoSection{\sigma}\bigr)^{-1}(0)\cap\ProtoBody{\sigma},
    \]
    and the exact-shell transport diffeomorphism
    \[
    \ProtoProj{\sigma}\big|_{\ProtoShell{\sigma}}
    :
    \ProtoShell{\sigma}
    \xrightarrow{\cong}
    \GermShell{\sigma};
    \]
    \item benchmark faithfulness, meaning preservation of the benchmark outputs \eqref{eq:fixedoutputs}, no second operative metric, no enlarged low-energy matter law, no change to the ordinary matter route, and no premature promotion from adoption to completed derivation.
\end{enumerate}
\end{definition}

\begin{definition}[Recorded current germ-anchored exact-shell transport-section core]
\label{def:sectioncore}
Fix \(\sigma\in\{\Car,\Cur,\Hom\}\) and the recorded transport core of Definition \ref{def:core}. The recorded current germ-anchored exact-shell transport-section core \(\ProtoSectionCore\) for sector \(\sigma\) consists of:
\begin{enumerate}
    \item the canonical exact-shell transport section
    \[
    \ProtoTransportSec{\sigma}
    \equiv
    \left(
    \ProtoProj{\sigma}\big|_{\ProtoShell{\sigma}}
    \right)^{-1}
    :
    \GermShell{\sigma}\to\ProtoShell{\sigma}\subset\ProtoBody{\sigma}\subset\ProtoState{\sigma};
    \]
    \item benchmark faithfulness in the sense of Definition \ref{def:core}.
\end{enumerate}
The exact proto shell and the restricted shell projection back to \(\GermShell{\sigma}\) are understood as derived data rather than as separately primitive section-core data.
\end{definition}

\begin{remark}
Definitions \ref{def:fixed}--\ref{def:sectioncore} are fixed inputs. The primitive continuation data, the one operative metric, the ordinary matter route, the downstream seed package, the recorded current germ nucleus, the recorded current transport-core package, and the already-generated section core are not reopened here as primitive stopping points. The present note asks only whether the current proto-section dynamics can themselves be generated from still more primitive substrate structure.
\end{remark}

\section{Recorded Proto-Section Target}

\begin{definition}[Recorded exact-shell-transport-section-generating substrate proto-section dynamics]
\label{def:protosection}
Fix \(\sigma\in\{\Car,\Cur,\Hom\}\). The recorded exact-shell-transport-section-generating substrate proto-section dynamics package \(\ProtoSectionPackage\) for sector \(\sigma\) consists of the following data.
\begin{enumerate}
    \item A finite-dimensional Euclidean proto-section state space \(\ProtoSectionState{\sigma}\), a compact convex proto-section body
    \[
    \ProtoSectionBody{\sigma}\subset\ProtoSectionState{\sigma}
    \]
    with nonempty interior, and an affine surjection
    \[
    \ProtoSectionProj{\sigma}:\ProtoSectionState{\sigma}\to\ProtoState{\sigma}
    \]
    satisfying
    \[
    \ProtoSectionProj{\sigma}\bigl(\ProtoSectionBody{\sigma}\bigr)=\ProtoBody{\sigma}.
    \]
    For each \(x\in\GermBody{\sigma}\), write
    \[
    \ProtoSectionFiber{\sigma}{x}
    \equiv
    \left\{
    y\in\ProtoSectionBody{\sigma}:
    \ProtoProj{\sigma}\bigl(\ProtoSectionProj{\sigma}(y)\bigr)=x
    \right\}.
    \]
    Each such fiber is required to be nonempty, compact, and convex.
    \item A nonnegative smooth fiber-selection functional
    \[
    \ProtoSectionFunctional{\sigma}:\ProtoSectionState{\sigma}\to[0,\infty)
    \]
    such that, for every \(x\in\GermBody{\sigma}\), the restriction of \(\ProtoSectionFunctional{\sigma}\) to \(\ProtoSectionFiber{\sigma}{x}\) is strictly convex and attains a unique minimizer. The resulting minimizer depends smoothly on \(x\).
    \item A closed embedded exact proto-section shell
    \[
    \ProtoSectionShell{\sigma}\subset\ProtoSectionBody{\sigma}
    \]
    satisfying
    \[
    \ProtoSectionShell{\sigma}
    =
    \left\{
    y\in\ProtoSectionBody{\sigma}:
    \ProtoProj{\sigma}\bigl(\ProtoSectionProj{\sigma}(y)\bigr)\in\GermShell{\sigma},
    \;
    \ProtoSectionFunctional{\sigma}(y)
    =
    \min_{z\in\ProtoSectionFiber{\sigma}{\ProtoProj{\sigma}(\ProtoSectionProj{\sigma}(y))}}
    \ProtoSectionFunctional{\sigma}(z)
    \right\},
    \]
    and
    \[
    \ProtoSectionProj{\sigma}\big|_{\ProtoSectionShell{\sigma}}
    :
    \ProtoSectionShell{\sigma}
    \xrightarrow{\cong}
    \ProtoShell{\sigma}
    \]
    is a diffeomorphism.
    \item Benchmark faithfulness, meaning preservation of the benchmark outputs \eqref{eq:fixedoutputs}, no second operative metric, no enlarged low-energy matter law, no change to the ordinary matter route, and no premature promotion from adoption to completed derivation.
\end{enumerate}
\end{definition}

\begin{remark}
Definition \ref{def:protosection} is the fixed target already carried by the current active-primary theorem. The present note does not spend its move on another theorem about the already-generated section core. It records the current proto-section package only because the deeper pre-proto-section dynamics are defined relative to it.
\end{remark}

\section{Proto-Section-Generating Substrate Pre-Proto-Section Dynamics}

\begin{definition}[Proto-section-generating substrate pre-proto-section dynamics]
\label{def:preprotosection}
Fix \(\sigma\in\{\Car,\Cur,\Hom\}\). A \emph{proto-section-generating substrate pre-proto-section dynamics package} \(\PreProtoSectionPackage\) for sector \(\sigma\) consists of the following data.
\begin{enumerate}
    \item A finite-dimensional Euclidean pre-proto-section state space \(\PreProtoSectionState{\sigma}\), a compact convex pre-proto-section body
    \[
    \PreProtoSectionBody{\sigma}\subset\PreProtoSectionState{\sigma}
    \]
    with nonempty interior, and an affine surjection
    \[
    \PreProtoSectionProj{\sigma}:\PreProtoSectionState{\sigma}\to\ProtoSectionState{\sigma}
    \]
    satisfying
    \[
    \PreProtoSectionProj{\sigma}\bigl(\PreProtoSectionBody{\sigma}\bigr)=\ProtoSectionBody{\sigma}.
    \]
    For each \(y\in\ProtoSectionBody{\sigma}\), write
    \[
    \PreProtoSectionFiber{\sigma}{y}
    \equiv
    \bigl(\PreProtoSectionProj{\sigma}\bigr)^{-1}(y)\cap\PreProtoSectionBody{\sigma}.
    \]
    Each such fiber is required to be nonempty, compact, and convex.
    \item A nonnegative smooth pre-proto-section functional
    \[
    \PreProtoSectionFunctional{\sigma}:\PreProtoSectionState{\sigma}\to[0,\infty)
    \]
    such that, for every \(y\in\ProtoSectionBody{\sigma}\), the restriction of \(\PreProtoSectionFunctional{\sigma}\) to \(\PreProtoSectionFiber{\sigma}{y}\) is strictly convex and attains a unique minimizer. The resulting minimizer depends smoothly on \(y\); write it as
    \[
    \PreProtoSectionProtoMin{\sigma}:\ProtoSectionBody{\sigma}\to\PreProtoSectionBody{\sigma}.
    \]
    These data satisfy
    \[
    \PreProtoSectionProj{\sigma}\circ\PreProtoSectionProtoMin{\sigma}
    =
    \mathrm{id}_{\ProtoSectionBody{\sigma}},
    \]
    and the effective minimum value recovers the recorded proto-section functional on the benchmark-relevant body:
    \[
    \ProtoSectionFunctional{\sigma}(y)
    =
    \PreProtoSectionFunctional{\sigma}\bigl(\PreProtoSectionProtoMin{\sigma}(y)\bigr)
    =
    \min_{u\in\PreProtoSectionFiber{\sigma}{y}}
    \PreProtoSectionFunctional{\sigma}(u)
    \qquad
    \text{for every }
    y\in\ProtoSectionBody{\sigma}.
    \]
    \item A closed embedded exact pre-proto-section shell
    \[
    \PreProtoSectionShell{\sigma}\subset\PreProtoSectionBody{\sigma}
    \]
    satisfying
    \[
    \PreProtoSectionShell{\sigma}
    =
    \left\{
    u\in\PreProtoSectionBody{\sigma}:
    \begin{array}{l}
    \PreProtoSectionProj{\sigma}(u)\in\ProtoSectionShell{\sigma},\\[0.2em]
    \PreProtoSectionFunctional{\sigma}(u)
    =
    \displaystyle\min_{v\in\PreProtoSectionFiber{\sigma}{\PreProtoSectionProj{\sigma}(u)}}
    \PreProtoSectionFunctional{\sigma}(v)
    \end{array}
    \right\},
    \]
    and
    \[
    \PreProtoSectionProj{\sigma}\big|_{\PreProtoSectionShell{\sigma}}
    :
    \PreProtoSectionShell{\sigma}
    \xrightarrow{\cong}
    \ProtoSectionShell{\sigma}
    \]
    is a diffeomorphism.
    \item Benchmark faithfulness, meaning preservation of the benchmark outputs \eqref{eq:fixedoutputs}, no second operative metric, no enlarged low-energy matter law, and presentation as a deeper substrate-side source for the current proto-section package rather than as a replacement benchmark theory.
\end{enumerate}
\end{definition}

\begin{remark}
The label \emph{pre-proto-section} is used here only as a local marker for a deeper layer below the current proto-section frontier. It does not by itself assert a completed microscopic ontology or a final primitive law. Its role is only to name the next sharper package below the current proto-section dynamics.
\end{remark}

\section{Induced Recovery of the Current Proto-Section Package}

\begin{proposition}[The effective proto-section functional is induced by fiberwise pre-proto-section minimization]
\label{prop:functional}
Fix a benchmark-faithful pre-proto-section package. Then the effective minimum-value functional
\[
y\longmapsto
\min_{u\in\PreProtoSectionFiber{\sigma}{y}}
\PreProtoSectionFunctional{\sigma}(u)
\]
is smooth on \(\ProtoSectionBody{\sigma}\) and agrees there with the recorded proto-section functional \(\ProtoSectionFunctional{\sigma}\).
\end{proposition}

\begin{proof}
Definition \ref{def:preprotosection}(2) states that the restriction of \(\PreProtoSectionFunctional{\sigma}\) to each fiber \(\PreProtoSectionFiber{\sigma}{y}\) has a unique minimizer and that this minimizer depends smoothly on \(y\). Hence the minimum-value map is smooth and equals
\[
\PreProtoSectionFunctional{\sigma}\bigl(\PreProtoSectionProtoMin{\sigma}(y)\bigr).
\]
The same definition requires this quantity to equal the recorded \(\ProtoSectionFunctional{\sigma}(y)\) for every \(y\in\ProtoSectionBody{\sigma}\).
\end{proof}

\begin{proposition}[Projected exact-shell transport recovers the recorded proto-section shell]
\label{prop:shell}
Fix a benchmark-faithful pre-proto-section package. Then
\[
\PreProtoSectionProj{\sigma}\bigl(\PreProtoSectionShell{\sigma}\bigr)
=
\ProtoSectionShell{\sigma},
\]
and
\[
\PreProtoSectionProj{\sigma}\big|_{\PreProtoSectionShell{\sigma}}
:
\PreProtoSectionShell{\sigma}\xrightarrow{\cong}\ProtoSectionShell{\sigma}
\]
is a diffeomorphism.
\end{proposition}

\begin{proof}
The displayed diffeomorphism is part of Definition \ref{def:preprotosection}(3). Its image is therefore exactly \(\ProtoSectionShell{\sigma}\).
\end{proof}

\begin{proposition}[Pre-proto-section dynamics canonically reproduce the current proto-section package]
\label{prop:package}
Fix a benchmark-faithful pre-proto-section package. Then the current exact-shell-transport-section-generating substrate proto-section dynamics of Definition \ref{def:protosection} are exactly the effective benchmark-relevant package induced on the recorded body \(\ProtoSectionBody{\sigma}\) by the deeper pre-proto-section dynamics. In particular:
\begin{enumerate}
    \item the recorded proto-section body and projection remain unchanged;
    \item the recorded proto-section functional on \(\ProtoSectionBody{\sigma}\) is the fiberwise minimum-value functional induced by the deeper pre-proto-section package;
    \item the recorded exact proto-section shell is the projected exact shell of the deeper pre-proto-section package.
\end{enumerate}
\end{proposition}

\begin{proof}
Item (1) is part of the fixed target data in Definition \ref{def:protosection} together with Definition \ref{def:preprotosection}(1). Item (2) is Proposition \ref{prop:functional}. Item (3) is Proposition \ref{prop:shell}. These statements are exactly the benchmark-relevant constituents of the current proto-section package.
\end{proof}

\begin{proposition}[The benchmark package remains fixed]
\label{prop:boundary}
For every benchmark-faithful pre-proto-section package, the operative metric remains exactly \(\CompletedMetric\), the ordinary matter route remains unchanged, no second operative metric is introduced, and the low-energy matter law is not enlarged.
\end{proposition}

\begin{proof}
This is part of the benchmark-faithfulness requirement in Definition \ref{def:preprotosection}(4). Equation \eqref{eq:fixedoutputs} remains unchanged on the full primitive forcing shell, so the deeper pre-proto-section package does not alter the benchmark exact-closure package.
\end{proof}

\begin{corollary}[The current section-core theorem applies unchanged]
\label{cor:section}
Fix a benchmark-faithful pre-proto-section package. Because Proposition \ref{prop:package} reproduces exactly the recorded proto-section package of Definition \ref{def:protosection}, the previously recorded theorem deriving the germ-anchored exact-shell transport-section core \(\ProtoSectionCore\) from exact-shell-transport-section-generating substrate proto-section dynamics applies unchanged. Consequently the current section core and all previously recorded downstream shell-transport consequences remain available unchanged.
\end{corollary}

\begin{proof}
The cited theorem requires exactly the proto-section package of Definition \ref{def:protosection} as its live input at current scope. Proposition \ref{prop:package} reproduces that package on the nose, so the same section-core recovery and downstream consequences follow without modification.
\end{proof}

\section{Main Theorem}

\begin{theorem}[Exact-shell-transport-section-generating substrate proto-section dynamics from proto-section-generating substrate pre-proto-section dynamics]
\label{thm:main}
Fix the primitive continuation data of Definition \ref{def:fixed}, the recorded downstream seed package of Definition \ref{def:seed}, the recorded current germ nucleus of Definition \ref{def:nucleus}, the recorded current transport core of Definition \ref{def:core}, and the recorded exact-shell-transport-section-generating substrate proto-section dynamics of Definition \ref{def:protosection}. Then, for each sector \(\sigma\in\{\Car,\Cur,\Hom\}\), every benchmark-faithful proto-section-generating substrate pre-proto-section dynamics package canonically produces the current proto-section package as the effective fiberwise minimum-value output of the deeper pre-proto-section layer. More precisely:
\begin{enumerate}
    \item the current exact-shell-transport-section-generating substrate proto-section dynamics are reproduced unchanged at benchmark scope by the deeper pre-proto-section dynamics;
    \item because the induced proto-section package is exactly the one already used by the current section-core theorem, the current germ-anchored exact-shell transport-section core \(\ProtoSectionCore\) and the previously recorded downstream consequences remain available unchanged;
    \item the same benchmark one-metric package remains fixed: the operative metric is still exactly \(\CompletedMetric\), the ordinary matter route is unchanged, there is no second operative metric, and the low-energy matter law is not enlarged;
    \item consequently the remaining honest burden on the active line is deeper again: derive or justify the proto-section-generating substrate pre-proto-section dynamics themselves from still more primitive substrate structure, or else prove a genuinely smaller theorem-level collapse, sharper necessity result, or sharper uniqueness result strictly inside that deeper package.
\end{enumerate}
\end{theorem}

\begin{proof}
Item (1) is Proposition \ref{prop:package}. Item (2) is Corollary \ref{cor:section}. Item (3) is Proposition \ref{prop:boundary}. Item (4) is the project-level consequence of the previous items together with the routing safeguard against counting another same-output deeper-origin relay as new front-line progress when the live benchmark-relevant datum is unchanged.
\end{proof}

\begin{corollary}[What remains open after this theorem]
\label{cor:next}
After Theorem \ref{thm:main}, the remaining honest burden is no longer to justify the exact-shell-transport-section-generating substrate proto-section dynamics themselves. Those dynamics are now the canonical effective output of a sharper pre-proto-section class. What remains open is why the proto-section-generating substrate pre-proto-section dynamics should hold, or whether an even sharper theorem can remove that still deeper pre-proto-section-level freedom through a smaller collapse, sharper necessity result, or sharper uniqueness result while preserving the same benchmark one-metric package and claim boundary.
\end{corollary}

\begin{proof}
Theorem \ref{thm:main} removes the current proto-section package as the unexplained stopping point on the active line. The next honest burden therefore lies strictly below it, unless a smaller theorem-level datum inside the pre-proto-section package can first be isolated and proved sufficient.
\end{proof}

\section{Conclusion}

The positive result of this note is narrower than a completed first-principles derivation and stronger than another same-output relay of the current section core. The current active-primary theorem had already shown that the germ-anchored exact-shell transport-section core is the canonical output of a deeper exact-shell-transport-section-generating substrate proto-section dynamics package. The present theorem then removes the next unexplained stopping point by showing that this current proto-section package itself is the effective projected output of a more primitive proto-section-generating substrate pre-proto-section dynamics.

The scope remains disciplined. The benchmark package is unchanged: the one operative metric is still exactly \(\CompletedMetric\), the ordinary matter route is unchanged, there is no second operative metric, and the low-energy matter law is not enlarged. Nothing here upgrades adoption to completed derivation or licenses stronger promotion language.

The open burden therefore moves one honest step further again. The next benchmark-facing manuscript should derive or justify the pre-proto-section dynamics from still more primitive substrate structure, unless the sharper honest gain available instead is a genuinely smaller theorem-level collapse, sharper necessity result, or sharper uniqueness result inside the observer-relevant pre-proto-section package. Until then, the present result should be read for what it is: a deeper proto-section-scope derivation on the active line of \TheoryName{}, not a claim that the full first-principles program is finished.

\end{document}
